\section{Phase Space Geometry}

\label{sec:geometry}

\subsection{Lagrangian Submanifolds and Equilibrium}

At market equilibrium, the state lies on a Lagrangian submanifold $\Lambda \subset \Gamma$ where $\bm{p} = \nabla S(\bm{q})$ for some generating function $S$. For HMM, the equilibrium manifold is:
\begin{equation}
\Lambda_{\text{eq}} = \{(\bm{q}, \bm{p}) : p_i = \kappa w_i / q_i, \; q_i = \bar{q}_i\}.
\end{equation}

Perturbations from $\Lambda_{\text{eq}}$ induce price oscillations whose frequency and damping are controlled by the parameters $(\lambda, \mu_i)$.

\subsection{KAM Stability}

The Kolmogorov--Arnold--Moser (KAM) theorem~\cite{arnold1963proof} guarantees that for sufficiently small perturbations from equilibrium, quasi-periodic orbits persist on invariant tori in phase space. In the HMM context, this means that moderate trading volume does not destroy the market's equilibrium structure.

\begin{theorem}[KAM Stability of HMM]
For the augmented Hamiltonian $\mathcal{H}_{\text{aug}}$ (\ref{eq:augmented}) with non-degenerate frequency vector $\bm{\omega} = (\omega_1, \ldots, \omega_n)$ satisfying the Diophantine condition $|\bm{k} \cdot \bm{\omega}| > \gamma |\bm{k}|^{-\tau}$ for all $\bm{k} \in \mathbb{Z}^n \setminus \{0\}$, invariant tori persist for perturbation strength $\epsilon < \epsilon_0(\gamma, \tau, n)$.
\end{theorem}

The Diophantine condition ensures that different resource types do not enter exact resonance, which would allow unbounded energy transfer between pools.

\subsection{Poincar\'e Sections and Trading Patterns}

We visualize market dynamics via Poincar\'e sections, projecting the high-dimensional phase space onto 2D slices. For a two-resource market (GPU, bandwidth), the section at $p_{\text{GPU}} = \text{const}$ reveals:

\begin{itemize}[leftmargin=1.1em]
    \item \textbf{Closed curves}: Stable, quasi-periodic trading around equilibrium.
    \item \textbf{Islands}: Resonant trading patterns (e.g., periodic batch jobs).
    \item \textbf{Chaotic regions}: High-frequency, unpredictable trading (typically thin).
\end{itemize}

The prevalence of closed curves in our simulations confirms that HMM dynamics are predominantly regular, ensuring predictable pricing behavior.

